% Part: applied-modal-logics
% Chapter: temporal-logic
% Section: possible-histories

\documentclass[../../../include/open-logic-section]{subfiles}

\begin{document}

\olfileid{aml}{tl}{poss}

\olsection{সম্ভাব্য ইতিহাস}

এতক্ষণ আমরা কালগত যুক্তিবিদ্যার যে সম্পর্কভিত্তিক মডেল ব্যবহার করেছি, তা অত্যন্ত নমনীয়: অভিগম্যতা-সম্পর্কের ওপর কোনো বিধিনিষেধ দিতে হয় না। ফলে কালগত মডেল অতীত ও ভবিষ্যৎ—উভয় দিকেই শাখাবিশিষ্ট হতে পারে। কিন্তু শাখাবিভাজনের আরও ``মোডাল'' ধারণা নিতে চাইতে পারি, যেখানে ঘটনাক্রমগুলিকে সম্ভাব্য ইতিহাস হিসেবে বিবেচনা করা হয়। এর জন্য ভাষা বদলানো আবশ্যক নয়; যদিও আমাদের ``সাধারণ'' মোডাল অপারেটর $\Box$ ও~$\Diamond$ যোগ করা যেতে পারে। আবার সময়ের সঙ্গে কর্তাদের জ্ঞান কীভাবে বদলায়, তা উপস্থাপনের জন্য জ্ঞানতাত্ত্বিক অভিগম্যতা-সম্পর্ক যোগ করার কথাও ভাবতে পারি।

\begin{defn}
  কালগত ভাষার একটি \emph{সম্ভাব্য ইতিহাসের মডেল} হল ত্রয়ী
  $\mModel{M} = \tuple{T, C, V}$, যেখানে
  \begin{enumerate}
    \item $T$ একটি অশূন্য সমষ্টি; এর সদস্যদের সময়ের অবস্থা হিসেবে ব্যাখ্যা করা হয়।
    \item $C$ একটি পদ্ধতির গণনাপথ, অর্থাৎ \emph{সম্ভাব্য ইতিহাস}-গুলির সমষ্টি। অন্যভাবে বললে, $C$~হল অবস্থা $s_1$, $s_2$,~$s_3$, \dots-এর ক্রম~$\sigma$-গুলির একটি সমষ্টি, যেখানে প্রত্যেক $s_i \in T$।
    \item $V$ এমন একটি অপেক্ষক, যা প্রত্যেক !!{propositional
      variable}~$p$-কে সময়ের বিন্দুগুলির একটি সমষ্টি~$V(p)$ নির্ধারণ করে দেয়।
  \end{enumerate}
  বিষয়টি সহজ রাখতে আমরা সাধারণত আরও ধরে নেব, কোনো ইতিহাস~$C$-তে থাকলে তার সব শেষাংশও সেই সমষ্টিতে থাকে। যেমন, $s_1$, $s_2$,~$s_3$ যদি~$C$-এর একটি ক্রম হয়, তবে $s_2$, $s_3$ এবং~$s_3$-ও তাতে আছে। আবার দুটি অবস্থা $s_i$ ও~$s_j$ কোনো ক্রম~$\sigma$-তে থাকলে, $i < j$ হলে আমরা লিখি $s_i \prec_\sigma s_j$। $t \in V(p)$ হলে বলি, $p$ বিন্দু~$t$-তে \emph{সত্য}।
\end{defn}

প্রাসঙ্গিক পরিবর্তন একটিই: মডেল~$\mModel M$-এর কোনো কালবিন্দু~$t$-তে !!a{formula}-এর সত্যতা নির্ণয় করি এমন একটি ইতিহাস~$\sigma$-এর সাপেক্ষে, যাতে $t$ একটি অবস্থা হিসেবে আছে। তবে !!{propositional
variable}s বা সত্য-অপেক্ষক সংযোজকের অর্থতত্ত্ব বদলাতে হয় না। এগুলি \olref[sem]{defn:tmodels}-এর মতোই থাকে, কারণ এগুলির কোনোটিতেই~$\sigma$-র উল্লেখ নেই। কিন্তু আমাদের ভবিষ্যৎ-অপারেটর~$\Ftemp$-কে এখন নতুন করে সংজ্ঞায়িত করি এবং এই ইতিহাসগুলির সাপেক্ষে~$\Diamond$ অপারেটর যোগ করি।

\begin{defn}\ollabel{defn:phmodels}
  $\mModel M$-তে \emph{$t, \sigma$-তে !!a{formula}~$!A$-এর সত্যতা}, প্রতীকে:
  $\mSat{M}{!A}[t, \sigma]$:
  \begin{enumerate}
  \item\ollabel{defn:sub:phmodels-f} \indcase{!A}{\Ftemp
    !B}{$\mSat{M}{\indfrm}[t, \sigma]$ তখনই এবং কেবল তখনই, যখন এমন কোনো $t' \in T$ আছে যাতে $t \prec_\sigma t'$ এবং $\mSat{M}{!B}[t', \sigma]$।}
  \item\ollabel{defn:sub:phmodels-diamond} \indcase{!A}{\Diamond
    !B}{$\mSat{M}{\indfrm}[t, \sigma]$ তখনই এবং কেবল তখনই, যখন এমন কোনো $\sigma' \in C$ আছে যাতে $t$ ঘটে এবং $\mSat{M}{!B}[t, \sigma']$।}
  \end{enumerate}
\end{defn}

অন্য কালগত ও মোডাল অপারেটরও একইভাবে সংজ্ঞায়িত করা যায়। তবে এখন কালরূপ ও মোডাল অবস্থা মিশিয়ে দাবি প্রকাশ করতে পারি। যেমন, ``$p$~ঘটবে না, কিন্তু ঘটতে পারত''—এটি $\lnot \Ftemp p \land \Diamond \Ftemp p$ সূত্র দিয়ে দেখানো যায়। যে বিন্দু ও ইতিহাসে পরবর্তী কোনো অবস্থায় $p$ সত্য হয় না, কিন্তু বিকল্প একটি ইতিহাসে পরে $p$ সত্য হবে, সেখানে এই সূত্র সত্য।

\end{document}
